\section{From resolutions to ambient right triviality}
\label{sec:ambient-lifting}

This section supplies the final topological step. Its input is complete
transport on the punctured zero fibers; its output is a trivialization of
the functions on ambient neighborhoods. We first make that passage on fixed
closed-ball tubes. We then construct the required transport from either
Whitney regularity or a simultaneous resolution, and finally remove a finite
base change. No estimate for a vector field at the singular section is needed.

\begin{lemma}[Ambient lifting on a closed-ball tube]
\label{lem:ambient-tube}
Let $G(x,u)$ be holomorphic on a neighborhood of
$B\times D$, where $B\subset\mathbb C^n$ is a closed ball centered at $0$
and $D\subset\mathbb C$ is a disk centered at $0$.
Suppose that $G(0,u)=0$ and that $d_xG(x,u)\ne0$ for $x\ne0$.
Assume that, for some $\eta>0$, the levels
$G(\cdot,u)=c$, $|c|<\eta$, meet $\partial B$ transversely.
Put
\[
 T_u=\{x\in B:|G(x,u)|<\eta\},\qquad
 Z^*=\{(x,u)\in(B\setminus\{0\})\times D:G(x,u)=0\}.
\]
Suppose there are smooth real vector fields $v_1,v_i$ on $Z^*$,
tangent to its sphere boundary, projecting to $1,i$ on $D$.
Suppose moreover that their transport along each segment $[0,u]$
and its reverse exists for the entire segment in $Z^*$.
More precisely, this is the transport of
$\operatorname{Re}(u)v_1+\operatorname{Im}(u)v_i$ for time one,
starting over $0$, and its reverse starting over $u$.
Then, after shrinking $D$ if necessary, there are homeomorphisms
\[
 H_u:T_0\longrightarrow T_u,
 \qquad G(H_u(x),u)=G(x,0),
\]
such that $H_0=\mathrm{id}$, $H_u(0)=0$, and both the family and its
inverse are jointly continuous. Restriction to the interiors of the tubes
is a parameter-preserving ambient homeomorphism near the origin section.
\end{lemma}

\begin{proof}
On the ambient tube with the origin section removed, $(G,u)$ is a
smooth submersion, including on the sphere boundary. Extend each $v_a$,
$a\in\{1,i\}$, to a smooth ambient field $V_a$ satisfying
\[
 du(V_a)=a,\qquad dG(V_a)=0,
 \qquad V_a\text{ tangent to the sphere boundary}.
\]
Here the prescribed extension along $Z^*$ can be made in submersion
coordinates, with boundary coordinates where necessary. Equivalently,
subtract a preliminary lift and extend the remaining section of the
kernel bundle of $d(G,u)$ from the closed submanifold $Z^*$ of the axis
complement. A partition of unity patches these extensions: the specified
projections, boundary tangency, and prescribed values on $Z^*$ are affine
linear conditions. There is no assertion that $V_a$ extends to the axis.

For $u\in D$, integrate
$\operatorname{Re}(u)V_1+\operatorname{Im}(u)V_i$ for time one,
starting at $(x,0)$. For $x\ne0$ on the zero level, existence for the
whole interval is the hypothesis. For $G(x,0)=c\ne0$, the trajectory
stays in the compact set
\[
 \{(y,s):y\in B,\ s\in[0,u],\ G(y,s)=c\},
\]
which misses the axis. Smooth continuation and boundary tangency give
existence for the whole interval here as well. Reverse transport gives
the inverse on $T_u\setminus\{0\}$. Both maps depend smoothly on their
arguments off the axis, also in charts at the sphere boundary, and preserve
the value of $G$.

Set $H_u(0)=0$. To prove joint continuity there, let $x_j\to0$ and
$u_j\to u$. A subsequential limit $y$ of $H_{u_j}(x_j)$ exists in $B$,
and function preservation gives $G(y,u)=0$. If $y\ne0$, smooth dependence
of reverse transport near $(y,u)$ would imply that $x_j$ tends to its
nonzero endpoint over $0$, a contradiction. Thus every such limit is $0$.
The same argument proves inverse continuity: if $y_j\in T_{u_j}$ tends
to $0$, any nonzero subsequential limit of $H_{u_j}^{-1}(y_j)$ contradicts
smooth forward transport. Consequently the maps and inverses are jointly
continuous everywhere. Reverse transport also proves that each $H_u$ is
onto exactly $T_u$. Finally, the tube interiors form open neighborhoods of
the respective origins and their union is an open neighborhood of the
origin section. This proves the last assertion.
\end{proof}

\begin{corollary}[Whitney regular zero fibers]
\label{cor:whitney-right}
Let $G(x,u)$, with $x\in\mathbb C^n$, be a holomorphic family with
$G(0,u)=0$ and no spatial critical points off the origin section near $(0,0)$. If the pair
\[
 \bigl(\{G=0\}\setminus\{x=0\},\ \{x=0\}\bigr)
\]
satisfies Whitney conditions \textup{(a)} and \textup{(b)}, then $G$ is
ambiently topologically right trivial near that section, with the central
map equal to the identity.
\end{corollary}

\begin{proof}
Write $\rho(x,u)=\|x\|^2$. We claim that $(u,\rho)$ is a real
submersion on the punctured zero locus sufficiently near $(0,0)$.
Otherwise choose points $(x_j,u_j)\to(0,0)$ witnessing failure.
After taking a subsequence, the unit spatial vectors $x_j/\|x_j\|$
converge to $v$, the real tangent spaces of the total zero locus converge
to $L$, and the real tangent spaces of its fibers converge to $Q$.
Whitney \textup{(a)} implies that $L$ contains the two real base
directions. The fiber tangent spaces have real dimension $2n-2$ and lie
in $\ker du$, so
\[
 Q=L\cap\ker du:
\]
the inclusion holds by passage to limits, and both sides have dimension
$2n-2$. Whitney \textup{(b)}, applied to the secants from $(0,u_j)$,
gives $(v,0)\in L$, hence $(v,0)\in Q$. On the other hand, failure of
submersivity says that $d\rho$ vanishes on each fiber tangent space.
It follows that $(v,0)$ is orthogonal to $Q$, a contradiction to
$\|v\|=1$.

Take a sufficiently small closed ball and parameter disk inside this
submersion neighborhood. Smooth local right inverses to $d(u,\rho)$,
patched by a partition of unity, give lifts $v_1,v_i$ of the parameter
directions with $d\rho(v_a)=0$. Their trajectories preserve the positive
spatial radius. Along any fixed parameter segment they therefore remain
in a compact set disjoint from the axis, so transport is complete in both
directions. These fields are tangent to the sphere boundary. The zero
levels are transverse to that sphere; by compactness the same holds for
all sufficiently small nonzero levels and parameters, after shrinking.
Lemma~\ref{lem:ambient-tube} applies.
\end{proof}

\begin{proposition}[Lifts supplied by a resolution]
\label{prop:resolution-right}
Let $G(x,u)$ satisfy the analytic and transversality hypotheses of
Lemma~\ref{lem:ambient-tube}. Suppose that the zero family admits a
proper holomorphic map $b:M\to\{G=0\}$ over the parameter disk such that
\begin{enumerate}
\item $q:M\to D$ is a smooth family;
\item $b$ is an isomorphism away from the origin section;
\item $E=b^{-1}(\{x=0\})$ is a relative normal-crossing divisor.
\end{enumerate}
Here the last condition means that, in local coordinates consisting of
fiber coordinates and $u$, the support of $E$ is a union of coordinate
hyperplanes in the fiber coordinates. Then $G$ has the trivialization
of Lemma~\ref{lem:ambient-tube}.
\end{proposition}

\begin{proof}
Choose the ball and a smaller parameter disk within representatives of
$b$. Above the closed ball, $q$ is a smooth submersion with boundary:
near the sphere, $b$ is an isomorphism and the fibers meet the sphere
transversely. Relative normal-crossing coordinates give local lifts of
the two real parameter directions tangent to every branch of $E$.
Away from $E$ use submersion coordinates, and near the sphere use boundary
submersion coordinates. A partition of unity gives smooth lifts on $M$
tangent both to $E$ and to the sphere boundary.

Properness makes the inverse image of the closed ball over each compact
parameter segment compact. The smooth lifts are therefore complete along
that segment in both directions. Their flows preserve $E$ and its
complement: tangency preserves each local branch, and uniqueness applied
in both time directions prevents a trajectory in the complement from
reaching $E$. Push the fields down through the isomorphism off $E$.
These are precisely the complete fields on the punctured zero family
required by Lemma~\ref{lem:ambient-tube}.
\end{proof}

\begin{lemma}[Removal of a finite disk base change]
\label{lem:remove-base-change}
Let $F(x,t)$ be a holomorphic family with $F(0,t)=0$.
Suppose that, after the base change $t=u^e$, $e\ge1$, the family
$G(x,u)=F(x,u^e)$ has the tube trivialization $H_u$ of
Lemma~\ref{lem:ambient-tube}. Then $F$ is ambiently topologically right
trivial over the original disk, by a section-preserving homeomorphism
whose central map is the identity.
\end{lemma}

\begin{proof}
Put $\zeta=\exp(2\pi i/e)$. For real $s\ge0$ sufficiently small,
\[
 J_s=H_{\zeta s}^{-1}\circ H_s:T_0\longrightarrow T_0
\]
is a homeomorphism preserving $f_0=F(\cdot,0)$, and $J_0=\mathrm{id}$.
These compositions have exact common domains: the two target tubes
$T_s$ and $T_{\zeta s}$ coincide since $s^e=(\zeta s)^e$.
The families $J_s$ and $J_s^{-1}=H_s^{-1}\circ H_{\zeta s}$ are jointly
continuous.

For $t=r^e\exp(i\theta)$, $r>0$ and $0\le\theta\le2\pi$, define
\[
 \phi_t
 =H_{r\exp(i\theta/e)}\circ J_{r\theta/(2\pi)}.
\]
The two values at the angular seam agree, because
\[
 H_{\zeta r}\circ J_r=H_r.
\]
Figure~\ref{fig:transport-seam} displays this endpoint identity.
\begin{figure}[ht]
\centering
\begin{minipage}[c]{0.49\textwidth}
\centering
\setlength{\unitlength}{1pt}
\begin{picture}(200,120)(-5,-16)
\put(10,10){\line(1,0){75}}
\put(10,10){\line(1,2){33.5}}
\qbezier(85,10)(85,56)(43.5,77)
\qbezier[12](10,10)(42.5,28.75)(75,47.5)
\qbezier(30,10)(30,16)(27.3,20)
\put(10,10){\circle*{3}}
\put(85,10){\circle*{3}}
\put(43.5,77){\circle*{3}}
\put(75,47.5){\circle*{3}}
\put(4,3){\makebox(0,0){$0$}}
\put(92,10){\makebox(0,0)[l]{$r$}}
\put(48,82){\makebox(0,0)[l]{$\zeta r$}}
\put(85,48){\makebox(0,0)[l]{$u=re^{i\theta/e}$}}
\put(36,21){\makebox(0,0)[l]{$\theta/e$}}
\put(75,-9){\makebox(0,0){$r,\zeta r\ \longmapsto\ t=r^e$}}
\end{picture}
\end{minipage}\hfill
\begin{minipage}[c]{0.49\textwidth}
\[
\begin{CD}
 T_0 @>{J_r}>> T_0\\
 @V{H_r}VV @VV{H_{\zeta r}}V\\
 T_r @= T_{\zeta r}
\end{CD}
\]
\end{minipage}
\[
 \phi_{r^e e^{i\theta}}
   =H_{r e^{i\theta/e}}\circ J_{r\theta/(2\pi)},
 \qquad 0\leq\theta\leq2\pi.
\]
\caption{Removing the angular seam for $t=u^e$ (the sector is schematic, with $e>1$).
One turn in the $t$-disk lifts from $r$ to $\zeta r$ in the $u$-disk.
These endpoints have the same target tube.  The square records the endpoint
correction $J_r=H_{\zeta r}^{-1}H_r$; the path $s\mapsto J_s$ from
$J_0=\mathrm{id}$ interpolates that correction as $\theta$ runs from
$0$ to $2\pi$.  The two endpoint formulas for $\phi_t$ therefore agree.}
\label{fig:transport-seam}
\end{figure}

The inverse is
\[
 \phi_t^{-1}
 =J_{r\theta/(2\pi)}^{-1}\circ H_{r\exp(i\theta/e)}^{-1},
\]
and its seam identity is
$J_r^{-1}\circ H_{\zeta r}^{-1}=H_r^{-1}$.
Thus both maps glue. At $r=0$ every constituent map is the identity,
independently of $\theta$; joint continuity of the constituent families
and their inverses therefore gives joint continuity on the polar disk
with its central circle collapsed. Set $\phi_0=\mathrm{id}$.

For every $t$, the source is the same tube $T_0$ and the target is
$\{x\in B:|f_t(x)|<\eta\}$. Moreover,
\[
 f_t\circ\phi_t=f_0\circ J_{r\theta/(2\pi)}=f_0,
 \qquad \phi_t(0)=0.
\]
Restricting to the interiors gives the asserted homeomorphism germ
$(x,t)\mapsto(\phi_t(x),t)$ over the original disk.
\end{proof}

\begin{remark}
A finite map between smooth pointed disks has the form $t=u^e$ after
a coordinate change: write it as $u^e$ times a nonvanishing holomorphic
unit and absorb a holomorphic $e$th root of that unit into $u$.
Accordingly, the preceding lemma applies to a normalized disk branch of
the finite base change used for a simultaneous resolution. It requires
no deck-equivariance of that resolution or of its chosen lifts.
\end{remark}
